% =====================================================================================================
% Section II. What Scaling Data Identify
% Writer: section "identification" (2026-09-24). Sources: output/memos/m7_theory.md (+ review),
% paper/notes/theory_main_text.md, paper/sections/appendix_proofs.tex, output/memos/m6_montecarlo.md (+ review).
% Numbers verified against output/tables/m6_montecarlo_designA_*.csv and output/tables/m7_theory_*.csv;
% see paper/notes/writer_identification.md for the ledger.
% =====================================================================================================

\providecommand{\E}{\mathbb{E}}
\providecommand{\Var}{\operatorname{Var}}
\providecommand{\Cov}{\operatorname{Cov}}

\section{What Scaling Data Identify}\label{sec:ident}

Section~\ref{sec:framework} described the technology through its dual objects: the path slope $a$, the frontier
elasticity $\gamma$, the elasticity of substitution $\sigma^*$ and the compute-optimal ratio $M^*(C)$. This section
asks which of them a given body of data identifies. The answer turns on a feature that sets scaling data apart from
most production data: the inputs are chosen by agents who minimize a known cost, $C=6ND$, and many come close to
succeeding. Optimization removes the variation that identifies curvature (subsection A), at a calculable rate
(subsection B). The sign of biased technical change is nonetheless identified (subsection C), the biases of cross-lab
regressions can be signed (subsection D), and revealed inference demand can be bounded when the technology must be
extrapolated (subsection E).

We state results for a family that nests the Chinchilla and Kaplan laws,
\begin{equation}\label{eq:ident-family}
L \;=\; E + e^{-\omega}\Big[A\big(e^{\psi_N}N\big)^{-a_1} + B\big(e^{\psi_D}D\big)^{-b_1}\Big]^{\kappa}e^{\epsilon},
\end{equation}
where $\omega$ is Hicks-neutral productivity, $(\psi_N,\psi_D)$ is factor-augmenting productivity, both known to the lab
when it chooses inputs, and $\epsilon$ is evaluation and seed noise, which is not. Chinchilla is $\kappa=1$ with
$(a_1,b_1)=(\alpha,\beta)$, and the joint law of \citet{kaplan2020scaling} is $\kappa\approx0.103$. Write
$S\equiv a_1+b_1$, $a\equiv b_1/S$, $b\equiv1-a$, $\gamma\equiv\kappa a_1b_1/S$ and $\sigma^*=2/(2+S)$; with $\kappa=1$
these are the objects of Section~\ref{sec:framework}. Proofs are in Online Appendix A, whose results are keyed to the propositions below. Every statement was also checked
symbolically and by a numerical method that does not use the closed form; the replication package reports 61 checks.

\subsection{Functional Dependence from Optimizing Behavior}\label{sec:ident-fdep}

The technology has a simple shape in log inputs. With the transverse coordinate $\tau\equiv a_1n-b_1d$,
\begin{equation}\label{eq:ident-geometry}
\ln R \;=\; \kappa\Big[-\tfrac12\,(a_1n+b_1d)+g(\tau)\Big],\qquad g''(\tau)\in\big(0,\tfrac14\big],
\end{equation}
where $R$ is reducible loss net of productivity and noise (Lemma~\ref{lem:geometry} in Online Appendix A). Reducible loss is
therefore exactly log-linear along every line with direction $(b_1,a_1)$, which is the direction of the compute-optimal
expansion path, and all curvature is transverse to it. A run's transverse position is its wedge: at the run's own compute,
$\ln w=(S/2)\ln(M/M^*(C))=(1/\sigma^*-1)\ln(M/M^*(C))$. Call a sample \emph{on-path} if its runs share one technology
with Hicks-neutral productivity ($\psi\equiv0$), each run minimizes training compute for its loss (no inference demand
and no binding constraint), and compute varies over an interval.

\begin{proposition}[Functional dependence]\label{prop:fdep}
The Hessian of $\ln R$ in $(n,d)$ has rank one, and its null vector is the direction of the expansion path. In an
on-path sample:
\begin{enumerate}
\item[(i)] $n_i=\ln G+a(c_i-\ln6)$ and $d_i=-\ln G+b(c_i-\ln6)$. The choice data reveal the path slope $a$ and intercept
$G$, that is, the observed $M^*(C)$, and nothing else.
\item[(ii)] $\E[y\mid c]=\gamma(c-\ln6)-\ln K+\E[\omega\mid c]$. The frontier elasticity $\gamma$ is identified if and
only if $\E[\omega\mid c]$ is known up to a constant, as in a single-lab sweep.
\item[(iii)] If $\E[\omega\mid c]$ is constant and $\kappa=1$, outcomes identify $E$, $\gamma$ and $K$ ($K$ under the
normalization $\E[e^{-\omega+\epsilon}]=1$), and hence $\alpha=\gamma/a$ and $\beta=\gamma/b$, but only at second order: the Jacobian of the on-path mean of $\ln L$ with respect
to $(E,\ln A,\ln B,\alpha,\beta)$ has rank three, and $\sigma^*$ is first-order identified if and only if
$\alpha=\beta$. The ratio $A/B$, and with it the technology's own compute-optimal ratio, is identified only by imposing
that the observed path is optimal.
\item[(iv)] If $\kappa$ is free, then for every $S>0$ the member of \eqref{eq:ident-family} with $a_1=(1-a)S$, $b_1=aS$,
$\kappa=\gamma/(a(1-a)S)$ and $(A,B)$ chosen to match $(G,K)$ generates the same joint distribution of $(n,d,y)$. Every
$\sigma^*\in(0,1)$ is consistent with the data; no $\sigma^*\ge1$ is.
\end{enumerate}
\end{proposition}

A ray has no curvature. Labs that minimize cost place their runs on the expansion path, where parameters and tokens are
deterministic functions of compute, so the curvature of the isoquant, which $\sigma^*$ measures, is never observed. This
is the functional-dependence problem of \citet{ackerberg2015identification}: an input that is a deterministic function
of the conditioning variables has no identified coefficient without variation that moves it and nothing else.
\citet{kricheli2026tokens} show that fixed tokens-per-parameter designs make scaling-law fits ill-conditioned. Our
point is that the collinearity is behavioral: cost-minimizing labs generate such designs on their own, and the better
they optimize, the closer they come to it. The remedies are those of the IO literature. An IsoFLOP sweep, which varies
$N$ at fixed $C$, is an experimenter-made shock to the relative price of the two inputs. In observational data the
analogous shocks are wedges that move the mix without moving productivity: heterogeneous inference demand
(Proposition~\ref{prop:wedge}), data and memory constraints, and belief errors such as allocation by the Kaplan law.
Each is an exclusion restriction and must be defended as one.

Parts (iii) and (iv) organize the published record. Chinchilla's Approaches 1 and 2 recover $(a,G)$ from choices and
Approach 3 recovers $(E,A,B,\alpha,\beta)$ from outcomes, linked by $G^{\alpha+\beta}=\alpha A/(\beta B)$
(Proposition~\ref{prop:duality}). Outcomes pin down the exponents (standard
errors of 0.02 in \citealp{besiroglu2024chinchilla}) but not the levels $A$ and $B$ (standard errors of 125 and 1,293),
and the compute-optimal ratio inherits the fragility of $A/B$. That ratio is what the revealed-demand calculation needs,
and on the path it is identified only by assuming that the observed allocation is optimal, the scaling analog of the
first-order-condition approach of \citet{gandhi2020identification}. Without $\kappa=1$, part (iv) says, curvature is not
identified at all. Figure~\ref{fig:geometry}, panel (b), draws three equivalent members ($\sigma^*=0.60$, 0.74 and
0.85): they coincide on the path and diverge off it. At five times the compute-optimal ratio, reducible loss lies
1.4 percent above the common frontier for a member with $\sigma^*=0.89$, 4.1 percent for Chinchilla ($\sigma^*=0.74$)
and 18.2 percent for a member with $\sigma^*=0.33$. Section~\ref{sec:tech:chinchilla} shows that the Chinchilla data
display the same pattern: with $\kappa$ free, the likelihood is flat in $\sigma^*$ on runs near the path.

\subsection{Information: The Better Labs Optimize, the Less Their Data Reveal}\label{sec:ident-info}

Real sweeps and real labs deviate from the path. How much do the deviations reveal? Because the frontier is a power
law, the output of any run in a designed sweep ($\omega$ normalized to zero) decomposes exactly as
\begin{equation}\label{eq:ident-decomp}
y_i \;=\; \gamma\,(c_i-\ln6)-\ln K-\gamma\,\Phi_i-\epsilon_i,
\end{equation}
where $\Phi_i\equiv\ln(C_i/C_{\min}(R_i))$, the run's Farrell allocative loss, is the log of the compute it spends
relative to the cheapest way of reaching its loss. It depends on the run's wedge alone (Section~\ref{sec:framework}), and
$\Phi_i\approx\sigma^*(\ln w_i)^2/[4(1-\sigma^*)]$ for small wedges. The curvature parameter $S$, and
hence $\sigma^*$, enters output only through these losses.

\begin{proposition}[Information]\label{prop:info}
Let $y_i=F(n_i,d_i;\phi)-\epsilon_i$ with $\epsilon_i$ independent $N(0,V_\epsilon)$, $E$ known, and the family
\eqref{eq:ident-family} parameterized by $\phi=(S,a,\gamma,\ln G,\ln K)$.
\begin{enumerate}
\item[(i)] The scores for the location of the path and for curvature are
\begin{align*}
\frac{\partial y_i}{\partial\ln G}&=\frac{\gamma(1-w_i)}{b+aw_i}=-\gamma\ln w_i+O\big((\ln w_i)^2\big),\\
\frac{\partial y_i}{\partial S}&=-\frac{\gamma}{S}\,\Phi_i\,\big[1+O(|\ln w_i|)\big].
\end{align*}
\item[(ii)] To leading order, the efficient information for $\ln G$ is $\gamma^2/V_\epsilon$ times the residual sum of
squares of $\ln w_i$ after projection on the other scores. The efficient information for $S$ is at most
$\gamma^2\sum_i\Phi_i^2/(S^2V_\epsilon)$, which is $O\big(\sum_i(\ln w_i)^4\big)$.
\item[(iii)] If every run is on the path, both informations are zero. If the transverse deviations of a fixed design
are scaled by $v\to0$, the information about the path location is $O(v^2)$ and that about $\sigma^*$ is $O(v^4)$: the
standard error of $\ln\hat M^*$ grows like $v^{-1}$ and that of $\hat\sigma^*$ like $v^{-2}$.
\end{enumerate}
\end{proposition}

A run reveals where the expansion path lies to the extent that it is off the path, which is first order in its log wedge.
It reveals the path's curvature only to the extent that it wastes compute, and the waste is a Harberger triangle, second
order in the wedge. Information is the square of the score, so the rates are second and fourth order: halving the
dispersion of log wedges quarters the information about $M^*$ and divides that about $\sigma^*$ by sixteen. Labs with
common objectives that all train compute-optimally generate no information about either. This is the precise sense in
which the better labs optimize, the less their data reveal. The result makes the Ackerberg--Caves--Frazer escape route
quantitative: optimization errors or price shocks reveal the location of the path at the rate of their variance and its
curvature at the rate of their fourth moment. Designed sweeps make wedges large on purpose; the transverse spread of the
Chinchilla extraction is $\mathrm{sd}(\ln w)=0.48$ across its 240 runs. In a nine-budget, five-offset factorial design,
the numerical log-log slopes of the information in the transverse spread are 4.00 for $\sigma^*$ and 2.00 for
$\ln M^*$ at the design's central compute. The results are local and assume Gaussian noise and a known $E$.

\begin{figure}[tp]
\centering
\includegraphics[width=\textwidth]{m6_montecarlo_fig2_designs}
\caption{Identification by Design: Monte Carlo Evidence}\label{fig:designs}
\begin{figurenotes}[Notes]
Simulations of the \citet{besiroglu2024chinchilla} technology ($\sigma^*=0.737$, $M^*=18.1$ at $10^{24}$ FLOP) with
normal noise of standard deviation 0.0075 in $\ln L$. Every design has 90 runs on Chinchilla's nine IsoFLOP budgets
($6\times10^{18}$ to $3\times10^{21}$ FLOP) and the same total compute, $5.1\times10^{22}$ FLOP. Optimizing labs set
$\ln(D/N)$ at its compute-optimal value plus a normal error with standard deviation $s$ (on-path: $s=0.02$). IsoFLOP
$\pm16\times$ places ten model sizes log-spaced over $[N^*/16,16N^*]$ at each budget; the factorial grid crosses nine
values of $N$ with ten values of $D$; Kaplan-belief labs set $N\propto C^{0.73}$ (with $s=0.1$). Primal: Huber
($\delta=10^{-3}$) fit of the five-parameter law with $\kappa=1$, as in \citet{hoffmann2022training}. Dual: path slope
from the optima and frontier elasticity from $L^*(C)$, assuming optimality and $\kappa=1$. Panel A: RMSE of
$\hat\sigma^*$ (500 replications per design). Panel B: median profile likelihood-ratio statistic over $\sigma^*$ with
$\kappa$ free (150 replications); horizontal line, 95 percent $\chi^2_1$ critical value; dotted line, truth. Panel C:
RMSE of $\ln\hat M^*$ at $10^{24}$ FLOP; the dual estimator is exact at $s=0$ by construction. Panel D: 5th, 25th,
50th, 75th and 95th percentiles of the primal $\hat\sigma^*$. One of the primal estimator's nine starting values is the
truth, which flatters the near-path designs.
\end{figurenotes}
\begin{figurenotes}[Source]
Authors' simulations; design and full tables in Online Appendix C.
\end{figurenotes}
\end{figure}

Figure~\ref{fig:designs} shows what these rates mean at realistic noise. We simulate the technology of
\citet{besiroglu2024chinchilla} with log-loss noise calibrated to the Chinchilla extraction. Every design has 90 runs and
the same total compute; designs differ only in where runs sit on their isocost lines. With every run exactly
compute-optimal ($s=0$), the Huber fit of \citet{hoffmann2022training} ends on a corner, with one power-law term
switched off, in 51 percent of replications, and the RMSEs of $\hat\sigma^*$ and $\ln\hat M^*$ are 0.28 and 24
(panels A and C). On the path a single exponential in compute is reproduced by the truth and by either term alone, so
the failure is a ridge with boundary solutions, not mere imprecision. As optimizing labs make larger errors, the RMSE of
$\hat\sigma^*$ falls to 0.067, 0.025 and 0.015 at $s=0.05$, 0.1 and 0.2, and that of $\ln\hat M^*$ to 4.5 at $s=0.1$,
2.0 at $s=0.3$ (where the 5--95 percent range of $\hat M^*$ still spans a factor of about 850) and 0.59 at $s=1$. An
IsoFLOP design with the same compute gives 0.0036 and 0.16.

Two findings qualify the apparent precision of near-optimal data. First, for $0.1\le s\le0.3$ the precision of
$\hat\sigma^*$ (RMSE 0.013--0.025) is functional-form information. Under $\kappa=1$, $\alpha+\beta=\gamma/(a(1-a))$ and
$a(1-a)$ is flat near $a=\tfrac12$, so $\sigma^*\approx2/(2+4\gamma)$ is pinned down by the on-path frontier elasticity;
the dual estimator, which uses only $(a,\gamma)$, is as precise (RMSE about 0.012). With $\kappa$ free (panel B), the
95 percent profile confidence set for $\sigma^*$ is the entire grid $[0.50,0.95]$ in 88 percent of on-path replications
and 65 percent at $s=0.3$; it becomes informative at $s=1$ (median width 0.10), against widths of 0.016 and 0.019 for
the IsoFLOP and factorial designs. What counts as near-optimal depends on noise and design size: at $s=0.3$ the share
of flat profiles is 38, 65 and 90 percent at noise standard deviations of 0.005, 0.0075 and 0.015. Second, the dual
route, which reads the technology off the path by assuming optimality, is precise on the path (RMSE 0.011 for
$\sigma^*$ and 0.009 for $\ln M^*$) but biased when labs optimize against wrong beliefs: for labs that follow Kaplan's
rule, the biases are $+0.22$ in $a$, $-0.09$ in $\sigma^*$ and $-3.9$ in $\ln M^*$ (a fiftyfold understatement).

Standard inference fails on the path. The Wald covariance is singular, or the fit sits at a corner, in 35 percent of
on-path replications. A pairs bootstrap warm-started at each replication's estimate, the usual implementation, covers the
true $\sigma^*$ in 67 percent of them. Restarting every draw from the full set of starting values restores coverage
(92 percent in a check with 60 replications), but the median interval width is then 0.43, against 0.016 under the
IsoFLOP design. On the path, inference is either wrong or empty; with designed variation the Wald and bootstrap
intervals agree and cover 95--97 percent of the time. These problems afflict compute-optimal ladders and sweeps.
Released industry models are often far off the path (under the Chinchilla technology, Llama 3 8B is trained at roughly
100 times its compute-optimal ratio), and for them the difficulty is extrapolation (subsection E).

\subsection{Technical Change over Time}\label{sec:ident-dmr}

Time series of compute-optimal models are the natural data for measuring biased technical change. Let
$\mathcal B_t\equiv\partial\ln(F_n/F_d)/\partial t$ at fixed inputs denote the Hicks bias, the drift in the marginal rate
of technical substitution; $\mathcal B_t>0$ is parameter-using.

\begin{proposition}[Technical change]\label{prop:dmr}
Suppose compute-optimal models are observed at several dates.
\begin{enumerate}
\item[(i)] For any smooth technology $y=F(n,d,t)$ with interior compute optima,
\begin{equation}\label{eq:ident-drift}
\frac{\partial\ln M^*(C,t)}{\partial t}\;=\;-\,\frac{\sigma^*_t}{1-\sigma^*_t}\;\mathcal B_t,\qquad\sigma^*_t\in(0,1).
\end{equation}
The sign of the Hicks bias is identified by the drift of the compute-optimal tokens-per-parameter ratio at given compute.
Its magnitude is not identified without $\sigma^*_t$.
\item[(ii)] In family \eqref{eq:ident-family} with factor-augmenting progress $\psi_{N,t}=g_Nt$ and $\psi_{D,t}=g_Dt$,
the frontier drift $\gamma(g_N+g_D)$ and the path drift $ag_D-bg_N$, together with $a$ and $\gamma$, which are
identified at each date, identify $g_N$ and $g_D$ separately. The bias
$\mathcal B=S(ag_D-bg_N)$ and $\sigma^*$ are not identified: a family of technologies indexed by $S>0$ is
observationally equivalent at every date.
\item[(iii)] Progress is Hicks-neutral ($a_1g_N=b_1g_D$) if and only if $M^*(C)$ is constant over time. Neutral change is
a proportional shift of $A$ and $B$, not a shift in $E$.
\item[(iv)] The compute-equivalent gain between two technologies, each evaluated on its own frontier, is constant in
compute if and only if they share $E$ and $\gamma$; their individual exponents may differ.
\end{enumerate}
\end{proposition}

\citet[hereafter DMR]{diamond1978measurement} showed that time series of output and inputs cannot separate the
elasticity of substitution from the bias of technical change. Part (i) refines this for scaling. With constant relative
input prices, the drift of the chosen input ratio reveals the sign of the Hicks bias for any $\sigma>0$; what it cannot
reveal in general is which input is augmented faster, because that sign flips at $\sigma=1$. In scaling, FLOP
accounting fixes relative log-cost weights at one and interior optima force $\sigma^*<1$ (Lemma~\ref{lemma:sigma}), so
both signs are identified; only magnitudes involving $\sigma^*$ are not. Because parameters and data are gross
complements, data-augmenting progress raises the marginal product of parameters relative to tokens: it is
parameter-using and lowers the compute-optimal tokens per parameter. In the equivalent family of part (ii), members with
$\sigma^*$ from 0.33 to 0.89 share frontiers (to $10^{-14}$ in $\ln R^*$) and expansion paths at six dates and seven
budgets; with $g_N=0.25$ and $g_D=0.60$ per year, all have a parameter-using bias, of 0.05 to 0.74 per year. With
$\kappa=1$ the bias is $\beta g_D-\alpha g_N$; the textbook form $(1-1/\sigma)(g_N-g_D)$ applies only when
$\alpha=\beta$. Part (iv) says when a single compute-equivalent multiplier summarizes progress: the condition involves
$E$ and $\gamma$, not the individual exponents.

The sign result holds relative prices fixed. If the effective cost of tokens relative to parameters changes, through
data-acquisition costs, a tightening data constraint or inference demand, the drift of $M^*$ mixes price responses with
bias and the DMR problem returns. The drift must be measured on training-optimal allocations, such as IsoFLOP minima at
comparable compute across vintages, not on deployed models, and part (ii) needs several budgets at each date. Section~\ref{sec:tech} applies part (iii) to data quality, with corpora in place of dates, and
Section~\ref{sec:obs} shows the DMR ridge in the cross-lab data of \citet{ho2024algorithmic}.

\subsection{Observational Data: Transmission, Selection, and Proxies}\label{sec:ident-obs}

Cross-lab data add the classic endogeneity of \citet{marschak1944random}: labs know their productivity when they choose
compute. Let runs lie on the expansion path of a common technology with Hicks-neutral productivity, and let compute follow
$c_i=\pi_0+\pi_1\omega_i+\eta_i$, with $(\omega,\eta,\epsilon)$ mutually uncorrelated and variances
$(V_\omega,V_\eta,V_\epsilon)$.

\begin{proposition}[Observational data]\label{prop:transmission}
Under these assumptions:
\begin{enumerate}
\item[(i)] (Transmission.) The OLS slope of $y$ on $c$ converges to $\gamma+\pi_1V_\omega/(\pi_1^2V_\omega+V_\eta)$: the
bias has the sign of $\pi_1$, and the dispersion of residual productivity is understated. The forward slope and the
reverse-regression estimate $\Var(y)/\Cov(c,y)$ bracket $\gamma$ if and only if
$-(1+V_\epsilon/V_\omega)/\gamma\le\pi_1\le0$.
\item[(ii)] (Allocation.) The OLS slope of $n$ on $c$ equals $a$ for every $\pi_1$. With factor-biased productivity it
converges to $a+\Cov(\chi,c)/[(\alpha+\beta)\Var(c)]$, where $\chi\equiv\beta\psi_D-\alpha\psi_N$.
\item[(iii)] (Selection.) If runs are released only when $y_i\ge\bar y$ and compute is exogenous, the released-sample
slope is at most $\gamma$ for any distribution of $\omega-\epsilon$, and lies in $[0,\gamma]$ if that distribution is
log-concave. Under joint normality, selection on $y$ leaves the reverse regression unchanged, and with $\pi_1>0$ the net
bias is positive if and only if the retained share of outcome variance exceeds a threshold $r^*$.
\item[(iv)] (Proxies.) Given compute, $\ln M_i=\ln M^*(C_i)+[2/(\alpha+\beta)](\ln w_i-\chi_i)$. With a common, known
wedge the mix reveals $\chi_i$ exactly, but it does not depend on $\omega_i$: no allocation-based proxy exists for
Hicks-neutral productivity. With heterogeneous wedges the mix reveals only $\ln w_i-\chi_i$.
\end{enumerate}
\end{proposition}

The sign of the transmission bias depends on the compute regime. OLS is consistent under exogenous budgets
($\pi_1=0$), as in a single-lab sweep, and overstates returns to compute when funding flows to productive labs
($\pi_1>0$). When labs buy compute until they reach a common loss target ($\pi_1=-1/\gamma$), more productive labs need
less of it and the slope is attenuated toward zero; if every lab hits the target exactly, the data carry no information
about $\gamma$. When budgets are provisioned before productivity is realized and respond to last period's productivity,
which persists at rate $\rho_\omega$, the effective $\pi_1$ is that response times $\rho_\omega$: the bias is upward but
attenuated. Simulated slopes match the closed form to within $4\times10^{-4}$ (Online Appendix C). The regime also
governs the familiar bracket. With classical measurement error, forward and reverse regressions bound the truth
\citep{klepper1984consistent}; here they do so under targets or exogenous budgets but not under funding. In an
illustrative calibration with $\gamma=0.178$ and $\pi_1=2$, the forward slope is 0.248 and the reverse estimate 0.403. The confounding
that \citet{whitfill2025note} and \citet{konig2026validity} warn about in observational scaling thus has a sign once the
compute regime is specified.

Part (ii) locates the damage. Hicks-neutral productivity scales both terms of the loss, leaves the marginal rate of
technical substitution unchanged and so never tilts the mix at given compute: transmission biases the frontier
elasticity, not the allocation exponent. Part (iii) is truncation on the dependent variable, of which dropping the
highest-loss runs from a sweep \citep{besiroglu2024chinchilla} is a mild instance. Selection pulls the slope down while
funding pushes it up, and either can dominate: in the simulations, a funding rule gives a net bias of $+0.100$ without
selection, $-0.006$ when half the runs are withheld and $-0.046$ when three quarters are.

Part (iv) explains why the control-function estimators of \citet{olley1996dynamics} and \citet{levinsohn2003estimating}
do not carry over. They invert a demand function that is strictly monotone in productivity. The natural candidate here,
the input mix, is flat in the productivity that matters most for returns to compute. The mix does invert the factor bias
$\chi$, but because $(n,d)$ and $(c,\ln M)$ are one-to-one, a first-stage control function in $(c,\ln M)$ absorbs the
entire technology, the \citet{ackerberg2015identification} problem again. Identification must come from timing, through
dynamic-panel moments with lagged outcomes (Online Appendix A), or from variation within model families, whose members
share productivity. Section~\ref{sec:obs} examines how far either route goes with public data.

\subsection{Partial Identification of Revealed Inference Demand}\label{sec:ident-pi}

Proposition~\ref{prop:wedge} turns a model's input mix into its planned inference demand, $T=3D(w-1)$. The wedge must be
evaluated at the model's own $(N,D)$, which usually lies outside the support of any design used to estimate the
technology.

\begin{proposition}[Partial identification]\label{prop:pi}
Consider a model at $(N,D)$ outside the design support.
\begin{enumerate}
\item[(i)] (Parametric.) A confidence set of level $1-p$ for $\theta=(\ln(A/B),\alpha,\beta)$ projects into the set
$\{3D[w(N,D;\theta)-1]\}$, which covers $T$ with probability at least $1-p$. At given compute,
$\ln\hat w=(\hat S/2)[\ln M-\ln\hat M^*(C)]$, so $\Var(\ln\hat w)$ is exactly quadratic in $\ln M$: the band widens
linearly in the log-extrapolation distance, with slope $z_{1-p/2}\,\mathrm{sd}(\hat S)/2$, and near the design support it
is dominated by the uncertainty in $\ln M^*$.
\item[(ii)] (Nonparametric.) If $F_{nn}\le0$, $F_{dd}\le0$ and $F_{nd}\ge0$, as throughout family
\eqref{eq:ident-family}, then $w$ is nonincreasing in $n$ and nondecreasing in $d$. Hence $T\ge3D\,[w(n_b,d_b)-1]$
for every support point $(n_b,d_b)$ with more parameters and fewer tokens than the model, where $w(n_b,d_b)$ is
estimated on the support without extrapolation. These shape restrictions
imply no finite upper bound, and without $F_{nd}\ge0$ the lower bound can fail.
\end{enumerate}
\end{proposition}

In part (i) the compute-optimal ratio, not the curvature, dominates. In a pairs bootstrap of the Huber fit to the
Chinchilla extraction (300 draws), at a point the size of Llama 3 8B ($M=1{,}875$, $C=7.2\times10^{23}$ FLOP),
$\mathrm{sd}(\ln\hat w)=0.161$; uncertainty in $\ln M^*$ alone contributes 0.143 and uncertainty in $\alpha+\beta$
alone 0.047. The illustration is heavily extrapolated: the design reaches $M=341$ only at compute up to
$1.3\times10^{22}$ FLOP, so at $7.2\times10^{23}$ every ratio is a 55-fold extrapolation in compute as well as in $M$.
Part (ii) delivers little on existing support. The best Chinchilla run with more parameters and fewer tokens than
Llama 3 8B has $w_b=0.76<1$ (evaluated with the fitted law), so the bound implies only $T\ge0$; informative bounds
require over-trained runs at large $N$.

The dominant source of uncertainty in revealed demand is thus $M^*$, the object that Proposition~\ref{prop:fdep}
identifies only through optimality and Proposition~\ref{prop:info} learns only at the rate of the variance of log
wedges. Section~\ref{sec:wedge} therefore calibrates $M^*(C)$ within model families, applying the first-order-condition
route of subsection A to a family's own compute-optimal law or flagship; this removes the largest term and leaves only
$\sigma^*$ to be imported from designed sweeps. The remaining error has a known direction: \citet{sardana2024beyond}
find that fitted laws overstate the value of extra tokens at extreme ratios, which understates $w$ there, so
extrapolated inference demand is conservative.

Taken together, the results rank the objects of Section~\ref{sec:framework} by how robustly they are identified. The
path slope $a$ is identified from near-optimal data and is immune to Hicks-neutral transmission bias. The frontier
elasticity $\gamma$ is identified within a lab; across labs its bias has the sign of the response of compute to
productivity. The curvature
$\sigma^*$ requires transverse variation and is learned at fourth order in it. The compute-optimal ratio $M^*$ requires
either the assumption of optimality or transverse variation, and it dominates the uncertainty of every quantity built on
the wedge. Table~\ref{tab:data} therefore reports how much transverse variation each dataset contains.
